Musical audio is a temporal signal: a track unfolds over time, and its content is determined not only by the set of timbres and instruments, but also by how they follow, return, and transform over time. Modern neural models encode audio as sequences of high-dimensional frame-level embeddings. Yet most downstream methods aggregate this sequence into a single vector, for example by mean pooling, max pooling, or weighted averaging. Such aggregation inevitably removes temporal structure.

An alternative view is to treat a track as a \textbf{trajectory} in feature space: a geometric object whose shape may carry information about musical content. This view is motivated by topological data analysis (TDA), especially persistent homology, which extracts shape invariants from point clouds (loops, connected components, voids) and summarizes them by persistence diagrams, which are stable under small perturbations~\cite{cohen2007stability}.

\textbf{Central question of the work:} does the topological organization of neural audio embedding trajectories carry meaningful information about music, and if so, what kind of information?

It is important to distinguish two claims from the outset. \emph{The topology of embedding trajectories} is what is measured here: geometric invariants of paths in a particular feature space. \emph{The topology of music} would be a much stronger claim: that these invariants are properties of music itself rather than artifacts of a chosen model. This work studies the former and deliberately avoids claiming the latter.

The motivation for testing genre signal is concrete. If genres systematically differ in the density of structural repetitions (for example, pop and rock with verse--chorus forms versus jazz with improvisation), then their embedding trajectories might show different recurrent geometry, detectable through~$H_1$. This is a testable hypothesis, not a guaranteed outcome, and its failure in the final results is itself informative.

The key methodological choice is that the objective is \textbf{not to develop a new method}, but to \textbf{rigorously assess the limits of applicability of TDA} to audio embeddings: with explicit controls against self-deception, with effect sizes rather than reliance on $p$-values, and with an honest account of both positive and negative results. Negative results are treated as a central contribution, not as a failure.

The rest of the report is organized as follows. Section~\ref{sec:litreview} reviews related work on TDA, audio embeddings, and delay embeddings. Section~\ref{sec:hypotheses} formulates the hypotheses tested in this work. Section~\ref{sec:data} describes the datasets. Section~\ref{sec:methodology} details the full pipeline from embedding extraction through surrogate controls to classification and cycle analysis. Section~\ref{sec:results} presents the experimental results for each hypothesis. Section~\ref{sec:discussion} discusses the findings, methodological lessons, and limitations. Section~\ref{sec:conclusion} summarizes the conclusions, and Section~\ref{sec:future} outlines directions for future work.
